\documentclass[12pt, titlepage]{article}
\usepackage[letterpaper, portrait, margin=1in]{geometry}
\usepackage{booktabs}
\usepackage{tabularx}
\usepackage{ltablex}
\usepackage{longtable}
\usepackage{hyperref}
\hypersetup{
    colorlinks,
    citecolor=black,
    filecolor=black,
    linkcolor=red,
    urlcolor=blue
}
\usepackage{longtable}
\usepackage{colortbl}
\usepackage{graphicx}
\usepackage{placeins}
\usepackage{array}
\usepackage{float}
\usepackage{colortbl}
\usepackage{caption}
\usepackage{graphicx}
\graphicspath{{./images/}}
\captionsetup[table]{width=0.9\textwidth,position=top}

\usepackage{listings}
\lstset{
  basicstyle=\ttfamily, % Use monospaced font
  breaklines=true,      % Enable line breaks
}
\usepackage{csvsimple}

\title{Reflection and Traceability Report on \progname}

\author{\authname}
\date{\today}

\input{../Comments}
\input{../Common}

\begin{document}

\maketitle
\pagenumbering{roman}

\iffalse
\plt{Reflection is an important component of getting the full benefits from a
learning experience.  Besides the intrinsic benefits of reflection, this
document will be used to help the TAs grade how well your team responded to
feedback.  Therefore, traceability between Revision 0 and Revision 1 is and
important part of the reflection exercise.  In addition, several CEAB (Canadian
Engineering Accreditation Board) Learning Outcomes (LOs) will be assessed based
on your reflections.}
\fi

\tableofcontents

\listoftables %if appropriate

\newpage

\pagenumbering{arabic}

\section{Changes in Response to Feedback}

\iffalse
\plt{Summarize the changes made over the course of the project in response to
feedback from TAs, the instructor, teammates, other teams, the project
supervisor (if present), and from user testers.}

\plt{For those teams with an external supervisor, please highlight how the feedback 
from the supervisor shaped your project.  In particular, you should highlight the 
supervisor's response to your Rev 0 demonstration to them.}

\plt{Version control can make the summary relatively easy, if you used issues
and meaningful commits.  If you feedback is in an issue, and you responded in
the issue tracker, you can point to the issue as part of explaining your
changes.  If addressing the issue required changes to code or documentation, you
can point to the specific commit that made the changes.  Although the links are
helpful for the details, you should include a label for each item of feedback so
that the reader has an idea of what each item is about without the need to click
on everything to find out.}

\plt{If you were not organized with your commits, traceability between feedback
and commits will not be feasible to capture after the fact.  You will instead
need to spend time writing down a summary of the changes made in response to
each item of feedback.}

\plt{You should address EVERY item of feedback.  A table or itemized list is
recommended.  You should record every item of feedback, along with the source of
that feedback and the change you made in response to that feedback.  The
response can be a change to your documentation, code, or development process.
The response can also be the reason why no changes were made in response to the
feedback.  To make this information manageable, you will record the feedback and
response separately for each deliverable in the sections that follow.}

\plt{If the feedback is general or incomplete, the TA (or instructor) will not
be able to grade your response to feedback.  In that case your grade on this
document, and likely the Revision 1 versions of the other documents will be
low.} 
\fi

\subsection{Problem Statement, Goals, and Development Plan} %------------------------------------------------------------------------

See Table~\ref{tab:PSGDPFeedback} for a summary of all feedback implemented. A majority of issues raised
came from our peers, a majority of which were implemented as suggested. See the linked issues for more detail in 
both the suggestions, as well as the pull requests for the corresponding changes.

\begin{longtable}{|p{2cm}|p{4cm}|p{5cm}|p{4cm}|}
\caption{Feedback Responses and Traceability for Problem Statement, Goals, Development Plan}
\label{tab:PSGDPFeedback} \\

\hline
\textbf{Feedback Source} & \textbf{Summary} & \textbf{Changes Made} & \textbf{Related Issues} \\ 
\hline
\endfirsthead

\hline
\multicolumn{4}{c}{\textit{Table \thetable\ continued}} \\
\hline
\textbf{Feedback Source} & \textbf{Summary} & \textbf{Changes Made} & \textbf{Related Issues} \\ 
\hline
\endhead



Peer & Enhanced high level goals & Revised all 5 goals to add justification, value \& stakeholder context & \href{https://github.com/OmarHassanAdelhamid/Five-of-a-Kind-capstone-project-/issues/37}{Issue \#37} \&
\href{https://github.com/OmarHassanAdelhamid/Five-of-a-Kind-capstone-project-/issues/44}{Issue \#44} \&
\href{https://github.com/OmarHassanAdelhamid/Five-of-a-Kind-capstone-project-/issues/45}{Issue \#45}
\newline\href{https://github.com/OmarHassanAdelhamid/Five-of-a-Kind-capstone-project-/pull/458/commits/054625d3b5fb38447b4884ebc11f88e6d637fe55}{Commit 054625d}
\newline\href{https://github.com/OmarHassanAdelhamid/Five-of-a-Kind-capstone-project-/pull/458}{Pull \#458} \\ \hline
Peer & Clarify label usage within the workflow plan  & Added an explanation for each label to clarify how it will be used in practice & \href{https://github.com/OmarHassanAdelhamid/Five-of-a-Kind-capstone-project-/issues/38}{Issue \#38} \&
\href{https://github.com/OmarHassanAdelhamid/Five-of-a-Kind-capstone-project-/issues/46}{Issue \#46} 
\newline\href{https://github.com/OmarHassanAdelhamid/Five-of-a-Kind-capstone-project-/pull/458/commits/488e2692100fe595d4640cdb2ee2132b56396931}{Commit 488e269}
\newline\href{https://github.com/OmarHassanAdelhamid/Five-of-a-Kind-capstone-project-/pull/458}{Pull \#458} \\ \hline
Peer & Risk mitigation strategies section & Separate existing mitigation strategies into a dedicated section \& strengthen the mitigation plan & \href{https://github.com/OmarHassanAdelhamid/Five-of-a-Kind-capstone-project-/issues/39}{Issue \#39}
\newline\href{https://github.com/OmarHassanAdelhamid/Five-of-a-Kind-capstone-project-/pull/458/commits/3f94db0e1390e4da2941a06e8a2ee3bd11095fd5}{Commit 3f94db0}
\newline\href{https://github.com/OmarHassanAdelhamid/Five-of-a-Kind-capstone-project-/pull/458}{Pull \#458} \\ \hline
Peer & Misalignment with problem statement \& POC plan & Refer to note in respective GitHub issue & \href{https://github.com/OmarHassanAdelhamid/Five-of-a-Kind-capstone-project-/issues/40}{Issue \#40}
\newline\href{https://github.com/OmarHassanAdelhamid/Five-of-a-Kind-capstone-project-/pull/458/commits/c9b03ee0ba0e7a01a8508e3a390b4bda4603bc95}{Commit c9b03ee}
\newline\href{https://github.com/OmarHassanAdelhamid/Five-of-a-Kind-capstone-project-/pull/458}{Pull \#458} \\ \hline
Peer & Justify technology choices & Justification was added for all selected technologies & \href{https://github.com/OmarHassanAdelhamid/Five-of-a-Kind-capstone-project-/issues/41}{Issue \#41}
\newline\href{https://github.com/OmarHassanAdelhamid/Five-of-a-Kind-capstone-project-/pull/458/commits/298906661109dc767dd2f90162c44bfa4a8b943e}{Commit 2989066}
\newline\href{https://github.com/OmarHassanAdelhamid/Five-of-a-Kind-capstone-project-/pull/458}{Pull \#458} \\ \hline
Peer & Team role management & Added section explaining how we will determine changes in team roles & \href{https://github.com/OmarHassanAdelhamid/Five-of-a-Kind-capstone-project-/issues/42}{Issue \#42}
\newline\href{https://github.com/OmarHassanAdelhamid/Five-of-a-Kind-capstone-project-/pull/458/commits/a59de22836fe40ed5557f470488e303978f7b7c3}{Commit a59de22}
\newline\href{https://github.com/OmarHassanAdelhamid/Five-of-a-Kind-capstone-project-/pull/458}{Pull \#458} \\ \hline
Peer & Depth in reflections & None (refer to note in respective GitHub issue for explanation) & \href{https://github.com/OmarHassanAdelhamid/Five-of-a-Kind-capstone-project-/issues/43}{Issue \#43}
\newline\href{https://github.com/OmarHassanAdelhamid/Five-of-a-Kind-capstone-project-/pull/458}{Pull \#458} \\ \hline
Peer & Reinforcement of coding standards & Add section explaining how coding standards will be enforced in team workflow & \href{https://github.com/OmarHassanAdelhamid/Five-of-a-Kind-capstone-project-/issues/47}{Issue \#47}
\newline\href{https://github.com/OmarHassanAdelhamid/Five-of-a-Kind-capstone-project-/pull/458/commits/6a21fdd6bfe9a5b52321a2a1a183a7c353f5511d}{Commit 6a21fdd}
\newline\href{https://github.com/OmarHassanAdelhamid/Five-of-a-Kind-capstone-project-/pull/458}{Pull \#458} \\ \hline
Peer & Contingency in Scheduling \& Deliverables & None (refer to note in respective GitHub issue for explanation) & \href{https://github.com/OmarHassanAdelhamid/Five-of-a-Kind-capstone-project-/issues/48}{Issue \#48}
\newline\href{https://github.com/OmarHassanAdelhamid/Five-of-a-Kind-capstone-project-/pull/458}{Pull \#458} \\ \hline


Team & Reference material assignment throughout & Rephrased throughout document to reference materials & \href{https://github.com/OmarHassanAdelhamid/Five-of-a-Kind-capstone-project-/issues/200}{Issue \#200}
\newline\href{https://github.com/OmarHassanAdelhamid/Five-of-a-Kind-capstone-project-/pull/279/commits/9afa9414904a5737602d6abb4abd919f6d082f43}{Commit 9afa941}
\newline\href{https://github.com/OmarHassanAdelhamid/Five-of-a-Kind-capstone-project-/pull/279}{Pull \#279} \\ \hline
TA & Add headliners to each section & Intro sentence added underneath each section & \href{https://github.com/OmarHassanAdelhamid/Five-of-a-Kind-capstone-project-/issues/202}{Issue \#202}
\newline\href{https://github.com/OmarHassanAdelhamid/Five-of-a-Kind-capstone-project-/pull/279/commits/9afa9414904a5737602d6abb4abd919f6d082f43}{Commit 9afa941}
\newline\href{https://github.com/OmarHassanAdelhamid/Five-of-a-Kind-capstone-project-/pull/279}{Pull \#279} \\ \hline
Team & Update POC Plan & Addressed in issue, pull request & \href{https://github.com/OmarHassanAdelhamid/Five-of-a-Kind-capstone-project-/issues/161}{Issue \#161}
\newline\href{https://github.com/OmarHassanAdelhamid/Five-of-a-Kind-capstone-project-/pull/252/commits/6e37f7bad766dfb49e4cdd5419fdb11361022e50}{Commit 6e37f7b}
\newline\href{https://github.com/OmarHassanAdelhamid/Five-of-a-Kind-capstone-project-/pull/252}{Pull \#252} \\ \hline
TA & Rephrase Goal 1 in Problem Statement & Split Goal 1 into two goals, importing and voxelization & \href{https://github.com/OmarHassanAdelhamid/Five-of-a-Kind-capstone-project-/issues/201}{Issue \#201}
\newline\href{https://github.com/OmarHassanAdelhamid/Five-of-a-Kind-capstone-project-/commit/c9c6c96d6cef5467090ec0b501c69eb57f151dff}{Commit c9c6c96}
\newline\href{https://github.com/OmarHassanAdelhamid/Five-of-a-Kind-capstone-project-/pull/444}{Pull \#444} \\ \hline
TA & Role assignment in the Team Members Role section in Development Plan & Assigned roles to team members & \href{https://github.com/OmarHassanAdelhamid/Five-of-a-Kind-capstone-project-/issues/203}{Issue \#203}
\newline\href{https://github.com/OmarHassanAdelhamid/Five-of-a-Kind-capstone-project-/pull/444/changes/bb14efbbfa486f5d750438d007afd9fe1afd357d}{Commit bb14efb}
\newline\href{https://github.com/OmarHassanAdelhamid/Five-of-a-Kind-capstone-project-/pull/444}{Pull \#444} \\ \hline
TA & Explain the process of branch naming in Development Plan & Added a section to the Development Plan to explain the process of branch naming, and an example of a branch name and how the workflow should work & \href{https://github.com/OmarHassanAdelhamid/Five-of-a-Kind-capstone-project-/issues/204}{Issue \#204}
\newline\href{https://github.com/OmarHassanAdelhamid/Five-of-a-Kind-capstone-project-/pull/444/changes/0268ad4cb4da45d1eed467240ab981dc53d6b717}{Commit 0268ad4}
\newline\href{https://github.com/OmarHassanAdelhamid/Five-of-a-Kind-capstone-project-/pull/444}{Pull \#444} \\ \hline
TA & Explain the GitHub Project columns in Development Plan & Added a section to the Development Plan to explain the GitHub Project columns and what each column is for & \href{https://github.com/OmarHassanAdelhamid/Five-of-a-Kind-capstone-project-/issues/205}{Issue \#205}
\newline\href{https://github.com/OmarHassanAdelhamid/Five-of-a-Kind-capstone-project-/pull/444/changes/a0ee521233ef5d3ab557145560fa44962be43c88}{Commit a0ee521}
\newline\href{https://github.com/OmarHassanAdelhamid/Five-of-a-Kind-capstone-project-/pull/444}{Pull \#444} \\ \hline
TA & Miscellaneous Formatting fixes & Made all the company names capital, and fixed quotations in latex & \href{https://github.com/OmarHassanAdelhamid/Five-of-a-Kind-capstone-project-/issues/206}{Issue \#206}
\newline\href{https://github.com/OmarHassanAdelhamid/Five-of-a-Kind-capstone-project-/pull/444/changes/cd1209d5c6c940b486f220eadd54d603ed884723}{Commit cd1209d}
\newline\href{https://github.com/OmarHassanAdelhamid/Five-of-a-Kind-capstone-project-/pull/444}{Pull \#444} \\ \hline
Peer & Task Assigning Strategy: Scheduling & Add description on how tasks will be decomposed and assigned to team members & \href{https://github.com/OmarHassanAdelhamid/Five-of-a-Kind-capstone-project-/issues/50}{Issue \#50}
\newline\href{https://github.com/OmarHassanAdelhamid/Five-of-a-Kind-capstone-project-/commit/3bb9f0c1fd7892a7597d2a2a6df1fcd8c329bef6}{Commit 3bb9f0c}
\newline\href{https://github.com/OmarHassanAdelhamid/Five-of-a-Kind-capstone-project-/pull/446}{Pull \#446} \\ \hline
Peer & Better external goals & Fix the external goals section to be more clear and concise. & \href{https://github.com/OmarHassanAdelhamid/Five-of-a-Kind-capstone-project-/issues/86}{Issue \#86}
\newline\href{https://github.com/OmarHassanAdelhamid/Five-of-a-Kind-capstone-project-/commit/7fc1aa47cb32687a8a36d497c8a638528c3abe0a}{Commit 7fc1aa4}
\newline\href{https://github.com/OmarHassanAdelhamid/Five-of-a-Kind-capstone-project-/pull/446}{Pull \#446} \\ \hline

Peer & Solution is introduced in Problem Statement & Remove the solution from the problem statement & \href{https://github.com/OmarHassanAdelhamid/Five-of-a-Kind-capstone-project-/issues/49}{Issue \#49}
\newline\href{https://github.com/OmarHassanAdelhamid/Five-of-a-Kind-capstone-project-/commit/110a1986c88a6d9deb9f61ec53206be98b980ce6}{Commit 110a198}
\newline\href{https://github.com/OmarHassanAdelhamid/Five-of-a-Kind-capstone-project-/pull/446}{Pull \#446} \\ \hline
\end{longtable}

\subsection{SRS and Hazard Analysis}

See Tables ~\ref{tab:SRSFeedback} \&~\ref{tab:HAFeedback} for a summary of all feedback implemented in the SRS and Hazard Analysis respectively. 
The SRS underwent radical changes, with many pieces of feedback resulting in the creation of new diagrams or significant 
rewording (eg. making requirements more verifiable). Many of these issues also cascaded into edits for the documents that 
followed. Significant formatting changes resulted as well to improve readability. The Hazard Analysis underwent similar changes, 
seeing the addition of new hazards and reformatting. See the linked issues and related pull requests for more detail.

\subsubsection{SRS}

\begin{longtable}{|p{2cm}|p{4cm}|p{5cm}|p{4cm}|} %------------------------------------------------------------------------
\caption{Feedback Responses and Traceability for SRS}
\label{tab:SRSFeedback} \\

\hline
\textbf{Feedback Source} & \textbf{Summary} & \textbf{Changes Made} & \textbf{Related Issues} \\ 
\hline
\endfirsthead

\hline
\multicolumn{4}{c}{\textit{Table \thetable\ continued}} \\
\hline
\textbf{Feedback Source} & \textbf{Summary} & \textbf{Changes Made} & \textbf{Related Issues} \\ 
\hline
\endhead

Team & Add missing general requirements & Addressed in issue, pull request & \href{https://github.com/OmarHassanAdelhamid/Five-of-a-Kind-capstone-project-/issues/160}{Issue \#160}
\newline\href{https://github.com/OmarHassanAdelhamid/Five-of-a-Kind-capstone-project-/pull/220/commits/4dcf99741fcdd6f39f03d536ecb05cabfae45d39}{Commit 4dcf997}
\newline\href{https://github.com/OmarHassanAdelhamid/Five-of-a-Kind-capstone-project-/pull/220}{Pull \#220} \\ \hline
Team & Add numbering to sections referenced in VnV Plan & Addressed in issue, pull request & \href{https://github.com/OmarHassanAdelhamid/Five-of-a-Kind-capstone-project-/issues/162}{Issue \#162}
\newline\href{https://github.com/OmarHassanAdelhamid/Five-of-a-Kind-capstone-project-/pull/220/commits/4dcf99741fcdd6f39f03d536ecb05cabfae45d39}{Commit 4dcf997}
\newline\href{https://github.com/OmarHassanAdelhamid/Five-of-a-Kind-capstone-project-/pull/220}{Pull \#220} \\ \hline
TA & Add more hyperlinks & Addressed in pull request & \href{https://github.com/OmarHassanAdelhamid/Five-of-a-Kind-capstone-project-/issues/208}{Issue \#208}
\newline\href{https://github.com/OmarHassanAdelhamid/Five-of-a-Kind-capstone-project-/pull/279/commits/6ab9071c341c1a43b016afde0d60a40277481fa6}{Commit 6ab9071}
\newline\href{https://github.com/OmarHassanAdelhamid/Five-of-a-Kind-capstone-project-/pull/279}{Pull \#279} \\ \hline
TA & Make traceability explicit & Added explicit traceability tables, addressed in pull request & \href{https://github.com/OmarHassanAdelhamid/Five-of-a-Kind-capstone-project-/issues/208}{Issue \#208}
\newline\href{https://github.com/OmarHassanAdelhamid/Five-of-a-Kind-capstone-project-/pull/279/commits/a99951220dd88428032f56d4d9d3810742d400b8}{Commit a999512}
\newline\href{https://github.com/OmarHassanAdelhamid/Five-of-a-Kind-capstone-project-/pull/279}{Pull \#279} \\ \hline
TA & Miscellaneous formatting fixes & Improved formatting, updated glossary; details in issue & \href{https://github.com/OmarHassanAdelhamid/Five-of-a-Kind-capstone-project-/issues/207}{Issue \#207}
\newline\href{https://github.com/OmarHassanAdelhamid/Five-of-a-Kind-capstone-project-/pull/279/commits/052ab2b02be49545426edea2cbbe768122050106}{Commit 052ab2b}
\newline\href{https://github.com/OmarHassanAdelhamid/Five-of-a-Kind-capstone-project-/pull/279}{Pull \#279} \\ \hline
TA, Peer & Add citations & Addressed in issue(s) & \href{https://github.com/OmarHassanAdelhamid/Five-of-a-Kind-capstone-project-/issues/212}{Issue \#212}
\newline\href{https://github.com/OmarHassanAdelhamid/Five-of-a-Kind-capstone-project-/issues/153}{Issue \#153}
\newline\href{https://github.com/OmarHassanAdelhamid/Five-of-a-Kind-capstone-project-/pull/279/commits/36222462e30c9a9a412a427c719f7115df62a3c2}{Commit 3622246}
\newline\href{https://github.com/OmarHassanAdelhamid/Five-of-a-Kind-capstone-project-/pull/279}{Pull \#279} \\ \hline
TA & Make low-level requirements verifiable & Reworded low-level requirements & \href{https://github.com/OmarHassanAdelhamid/Five-of-a-Kind-capstone-project-/issues/209}{Issue \#209}
\newline\href{https://github.com/OmarHassanAdelhamid/Five-of-a-Kind-capstone-project-/pull/279/commits/bbd8ff7e7ffe140aca18d3cd53480f9ecd4beeba}{Commit bbd8ff7}
\newline\href{https://github.com/OmarHassanAdelhamid/Five-of-a-Kind-capstone-project-/pull/279}{Pull \#279} \\ \hline
Peer & Make requirements verifiable & Reworded high-level requirements & \href{https://github.com/OmarHassanAdelhamid/Five-of-a-Kind-capstone-project-/issues/152}{Issue \#152}
\newline\href{https://github.com/OmarHassanAdelhamid/Five-of-a-Kind-capstone-project-/pull/279/commits/42effd17a34adde7426efea55106c82f3724c94e}{Commit 42effd1}
\newline\href{https://github.com/OmarHassanAdelhamid/Five-of-a-Kind-capstone-project-/pull/279}{Pull \#279} \\ \hline
Peer & Use Case Diagram Too detailed & Refactored diagram; addressed in pull request & \href{https://github.com/OmarHassanAdelhamid/Five-of-a-Kind-capstone-project-/issues/147}{Issue \#147}
\newline\href{https://github.com/OmarHassanAdelhamid/Five-of-a-Kind-capstone-project-/pull/279/commits/42effd17a34adde7426efea55106c82f3724c94e}{Commit 42effd1}
\newline\href{https://github.com/OmarHassanAdelhamid/Five-of-a-Kind-capstone-project-/pull/279}{Pull \#279} \\ \hline
TA, Peer & Expand societal benefits, goal model & Refactored relevant section; addressed in pull request & \href{https://github.com/OmarHassanAdelhamid/Five-of-a-Kind-capstone-project-/issues/214}{Issue \#214}
\newline\href{https://github.com/OmarHassanAdelhamid/Five-of-a-Kind-capstone-project-/issues/154}{Issue \#154}
\newline\href{https://github.com/OmarHassanAdelhamid/Five-of-a-Kind-capstone-project-/issues/155}{Issue \#155}
\newline\href{https://github.com/OmarHassanAdelhamid/Five-of-a-Kind-capstone-project-/pull/279/commits/836ae244e2b586d0e8fe8a698d4b32b65836123d}{Commit 836ae24}
\newline\href{https://github.com/OmarHassanAdelhamid/Five-of-a-Kind-capstone-project-/pull/279}{Pull \#279} \\ \hline
TA & Make system boundary clear & Added system boundary diagram & \href{https://github.com/OmarHassanAdelhamid/Five-of-a-Kind-capstone-project-/issues/213}{Issue \#213}
\newline\href{https://github.com/OmarHassanAdelhamid/Five-of-a-Kind-capstone-project-/pull/279/commits/741be088bb8af47f066f40b65198f5236b733b3a}{Commit 741be08}
\newline\href{https://github.com/OmarHassanAdelhamid/Five-of-a-Kind-capstone-project-/pull/279}{Pull \#279} \\ \hline
TA & UC Diagram clarity & Made arrow directions consistent, all actors no longer inherit from user; addressed in pull request. & \href{https://github.com/OmarHassanAdelhamid/Five-of-a-Kind-capstone-project-/issues/210}{Issue \#210}
\newline\href{https://github.com/OmarHassanAdelhamid/Five-of-a-Kind-capstone-project-/pull/442/commits/f5ab69ed317064499814f008bf446b8cde2b4949}{Commit f5ab69e}
\newline\href{https://github.com/OmarHassanAdelhamid/Five-of-a-Kind-capstone-project-/pull/442}{Pull \#442} \\ \hline
TA & SRS is missing formalization & Added a system state diagram to section S.3 & \href{https://github.com/OmarHassanAdelhamid/Five-of-a-Kind-capstone-project-/issues/215}{Issue \#215}
\newline\href{https://github.com/OmarHassanAdelhamid/Five-of-a-Kind-capstone-project-/pull/442/commits/30bb2fd418ad76cea34a67a6aaf95fdd11296b60}{Commit 30bb2fd}
\newline\href{https://github.com/OmarHassanAdelhamid/Five-of-a-Kind-capstone-project-/pull/442}{Pull \#442} \\ \hline
Peer & List required technologies & Closed with comment; SRS should be independent from the solution/implementation. & \href{https://github.com/OmarHassanAdelhamid/Five-of-a-Kind-capstone-project-/issues/151}{Issue \#151} \\ \hline
TA & Relate the Project Sprints to the User Stories & Added a section to the SRS that relates the project sprints to the user stories & \href{https://github.com/OmarHassanAdelhamid/Five-of-a-Kind-capstone-project-/issues/211}{Issue \#211} \newline\href{https://github.com/OmarHassanAdelhamid/Five-of-a-Kind-capstone-project-/pull/462/changes/96b81a428fece8bdf142340ca2d9e239212c390f}{Commit 96b81a4} \newline\href{https://github.com/OmarHassanAdelhamid/Five-of-a-Kind-capstone-project-/pull/462}{Pull \#462} \\ \hline

\end{longtable}

\subsubsection{Hazard Analysis}

\begin{longtable}{|p{2cm}|p{4cm}|p{5cm}|p{3cm}|} %------------------------------------------------------------------------
\caption{Feedback Responses and Traceability for Hazard Analysis}
\label{tab:HAFeedback} \\

\hline
\textbf{Feedback Source} & \textbf{Summary} & \textbf{Changes Made} & \textbf{Related Issues} \\ 
\hline
\endfirsthead

\hline
\multicolumn{4}{c}{\textit{Table \thetable\ continued}} \\
\hline
\textbf{Feedback Source} & \textbf{Summary} & \textbf{Changes Made} & \textbf{Related Issues} \\ 
\hline
\endhead

TA & Reformat tables, and add list of tables & Addressed in pull request & \href{https://github.com/OmarHassanAdelhamid/Five-of-a-Kind-capstone-project-/issues/216}{Issue \#216}
\newline\href{https://github.com/OmarHassanAdelhamid/Five-of-a-Kind-capstone-project-/pull/442/commits/94531dce8e16b328345c0290f14a1d5b71fdd965}{Commit 94531dc}
\newline\href{https://github.com/OmarHassanAdelhamid/Five-of-a-Kind-capstone-project-/pull/442}{Pull \#242} \\ \hline
TA & Fix roadmap timeline/priority inconsistencies & Reworked to match implementation, changed priorities accordingly & \href{https://github.com/OmarHassanAdelhamid/Five-of-a-Kind-capstone-project-/issues/217}{Issue \#217}
\newline\href{https://github.com/OmarHassanAdelhamid/Five-of-a-Kind-capstone-project-/pull/442/commits/94531dce8e16b328345c0290f14a1d5b71fdd965}{Commit 94531dc}
\newline\href{https://github.com/OmarHassanAdelhamid/Five-of-a-Kind-capstone-project-/pull/442}{Pull \#242} \\ \hline
Peer & Hazard: File external compatibility & Closed with comment; addressed in an existing hazard. & \href{https://github.com/OmarHassanAdelhamid/Five-of-a-Kind-capstone-project-/issues/149}{Issue \#149} \\ \hline
Peer & Roadmap Detail/Linking & Added hyperlinks & \href{https://github.com/OmarHassanAdelhamid/Five-of-a-Kind-capstone-project-/issues/158}{Issue \#158}
\newline\href{https://github.com/OmarHassanAdelhamid/Five-of-a-Kind-capstone-project-/pull/460/commits/9a17d63126d54891faf5438eb9b9f26e7b860020}{Commit 9a17d63}
\newline\href{https://github.com/OmarHassanAdelhamid/Five-of-a-Kind-capstone-project-/pull/460}{Pull \#460} \\ \hline
TA & Missing hazard - Incorrect properties & Added hazard and reworded security requirements to include & \href{https://github.com/OmarHassanAdelhamid/Five-of-a-Kind-capstone-project-/issues/219}{Issue \#219}
\newline\href{https://github.com/OmarHassanAdelhamid/Five-of-a-Kind-capstone-project-/pull/460/commits/ab0fe5610fd5f6aec7be4830c39b80b6081480b5}{Commit ab0fe56}
\newline\href{https://github.com/OmarHassanAdelhamid/Five-of-a-Kind-capstone-project-/pull/460}{Pull \#460} \\ \hline
TA & Unclear scope & Added high-level scope areas to further clarify & \href{https://github.com/OmarHassanAdelhamid/Five-of-a-Kind-capstone-project-/issues/218}{Issue \#218}
\newline\href{https://github.com/OmarHassanAdelhamid/Five-of-a-Kind-capstone-project-/pull/464/commits/7ed5f4a6ce8723e48c76929e7634bcc6f78f03f0}{Commit 7ed5f4a}
\newline\href{https://github.com/OmarHassanAdelhamid/Five-of-a-Kind-capstone-project-/pull/464}{Pull \#464} \\ \hline

\end{longtable}

\subsection{Design and Design Documentation}

See Table~\ref{tab:DDFeedback} for a summary of all feedback and changes for the Module Guide and Module
Interface Specification. Changes to these documents was quite extensive as this document received feedback twice;
with more time it is likely that even further changes and improvements would occur. Overall, changes include improving
consistency and alignment with the design, overarching changes to the design to improve modularity, reworking of modules
to align with MIS specifications, and formatting to improve readability. A handful of feedback was also addressed with comments,
which is explained in most detail in the linked issues where relevant as closing comments. Issues from Revision -1 were closed
with comments as some issues referred to modules or sections that were removed or entirely reworked.

\begin{longtable}{|p{2cm}|p{4cm}|p{5cm}|p{3cm}|} %------------------------------------------------------------------------
\caption{Feedback Responses and Traceability for Design Documentation}
\label{tab:DDFeedback} \\

\hline
\textbf{Feedback Source} & \textbf{Summary} & \textbf{Changes Made} & \textbf{Related Issues} \\ 
\hline
\endfirsthead

\hline
\multicolumn{4}{c}{\textit{Table \thetable\ continued}} \\
\hline
\textbf{Feedback Source} & \textbf{Summary} & \textbf{Changes Made} & \textbf{Related Issues} \\ 
\hline
\endhead

Peer & Broad Anticipated Change (AC9) & Minor reword, but left unchanged (see issue comment for full explanation) & \href{https://github.com/OmarHassanAdelhamid/Five-of-a-Kind-capstone-project-/issues/247}{Issue \#247}
\newline\href{https://github.com/OmarHassanAdelhamid/Five-of-a-Kind-capstone-project-/pull/303/commits/ec884b8971712fe48cbedb8852b482cfae0e8053}{Commit ec884b8}
\newline\href{https://github.com/OmarHassanAdelhamid/Five-of-a-Kind-capstone-project-/pull/303}{Pull \#303} \\ \hline
TA & Miscellaneous formatting changes & See issue for details; addressed in pull request & \href{https://github.com/OmarHassanAdelhamid/Five-of-a-Kind-capstone-project-/issues/274}{Issue \#274}
\newline\href{https://github.com/OmarHassanAdelhamid/Five-of-a-Kind-capstone-project-/pull/303/commits/a8a3263808c0330510317c935618e14e22cd44d8}{Commit a8a3263}
\newline\href{https://github.com/OmarHassanAdelhamid/Five-of-a-Kind-capstone-project-/pull/303}{Pull \#303} \\ \hline
Peer & Cross referencing between MG/MIS difficult & Added consistent numbering between the two & \href{https://github.com/OmarHassanAdelhamid/Five-of-a-Kind-capstone-project-/issues/251}{Issue \#251}
\newline\href{https://github.com/OmarHassanAdelhamid/Five-of-a-Kind-capstone-project-/pull/303/commits/a8a3263808c0330510317c935618e14e22cd44d8}{Commit a8a3263}
\newline\href{https://github.com/OmarHassanAdelhamid/Five-of-a-Kind-capstone-project-/pull/303}{Pull \#303} \\ \hline
Peer & Undefined data type & Removed references, addressed in commit/pull request & \href{https://github.com/OmarHassanAdelhamid/Five-of-a-Kind-capstone-project-/issues/248}{Issue \#248}
\newline\href{https://github.com/OmarHassanAdelhamid/Five-of-a-Kind-capstone-project-/pull/304}{Pull \#304}
\newline\href{https://github.com/OmarHassanAdelhamid/Five-of-a-Kind-capstone-project-/pull/310/commits/3024aea8910e0dafd3227d653f9d8359b277b5ec}{Commit 3024aea}
\newline\href{https://github.com/OmarHassanAdelhamid/Five-of-a-Kind-capstone-project-/pull/310}{Pull \#310} \\ \hline
TA & Event should be a datatype/enum & Addressed in commit & \href{https://github.com/OmarHassanAdelhamid/Five-of-a-Kind-capstone-project-/issues/277}{Issue \#277}
\newline\href{https://github.com/OmarHassanAdelhamid/Five-of-a-Kind-capstone-project-/pull/311/commits/3e21652aba486a639dda369b16c2b0d66f0f43e0}{Commit 3e21652}
\newline\href{https://github.com/OmarHassanAdelhamid/Five-of-a-Kind-capstone-project-/pull/311}{Pull \#311} \\ \hline
TA & Low level details (promise) & Removed all references, addressed in commit & \href{https://github.com/OmarHassanAdelhamid/Five-of-a-Kind-capstone-project-/issues/275}{Issue \#275}
\newline\href{https://github.com/OmarHassanAdelhamid/Five-of-a-Kind-capstone-project-/pull/310/commits/3024aea8910e0dafd3227d653f9d8359b277b5ec}{Commit 3024aea}
\newline\href{https://github.com/OmarHassanAdelhamid/Five-of-a-Kind-capstone-project-/pull/310}{Pull \#310} \\ \hline
TA & Hierarchy diagram issues (1) & Removed accidental loops, fixed error colors & \href{https://github.com/OmarHassanAdelhamid/Five-of-a-Kind-capstone-project-/issues/273}{Issue \#273}
\newline\href{https://github.com/OmarHassanAdelhamid/Five-of-a-Kind-capstone-project-/pull/310/commits/3024aea8910e0dafd3227d653f9d8359b277b5ec}{Commit 3024aea}
\newline\href{https://github.com/OmarHassanAdelhamid/Five-of-a-Kind-capstone-project-/pull/310}{Pull \#310} \\ \hline
Peer & Unclear distinction between modules & During rework of modules for Rev 0, these were entirely refactored & \href{https://github.com/OmarHassanAdelhamid/Five-of-a-Kind-capstone-project-/issues/246}{Issue \#246}
\newline\href{https://github.com/OmarHassanAdelhamid/Five-of-a-Kind-capstone-project-/issues/250}{Issue \#250 (duplicate)}
\newline\href{https://github.com/OmarHassanAdelhamid/Five-of-a-Kind-capstone-project-/pull/310/commits/3024aea8910e0dafd3227d653f9d8359b277b5ec}{Commit 3024aea}
\newline\href{https://github.com/OmarHassanAdelhamid/Five-of-a-Kind-capstone-project-/pull/310}{Pull \#310} \\ \hline
TA & Misidentified module & Module was merged into another during refactoring and re-identified; closed with comment & \href{https://github.com/OmarHassanAdelhamid/Five-of-a-Kind-capstone-project-/issues/271}{Issue \#271} \\ \hline
TA & Hierarchy diagram issues (2) & Organized to be a top-down DAG with fewer bends; is a vector graphic. & \href{https://github.com/OmarHassanAdelhamid/Five-of-a-Kind-capstone-project-/issues/427}{Issue \#427}
\newline\href{https://github.com/OmarHassanAdelhamid/Five-of-a-Kind-capstone-project-/pull/442/commits/6e37e20d778a381496a3eb90bab1a087da8fc3d6}{Commit 6e37e20}
\newline\href{https://github.com/OmarHassanAdelhamid/Five-of-a-Kind-capstone-project-/pull/442}{Pull \#442} \\ \hline
Team & Mistake in MG GraphicsAdapter description & Rewrote/replaced erroneous description. & \href{https://github.com/OmarHassanAdelhamid/Five-of-a-Kind-capstone-project-/issues/443}{Issue \#443}
\newline\href{https://github.com/OmarHassanAdelhamid/Five-of-a-Kind-capstone-project-/pull/442/commits/f5ab69ed317064499814f008bf446b8cde2b4949}{Commit f5ab69e}
\newline\href{https://github.com/OmarHassanAdelhamid/Five-of-a-Kind-capstone-project-/pull/442}{Pull \#442} \\ \hline
TA, Peer & Unclear/questionable unlikely changes in MG & Rewrote section; all new unlikely changes noted. & \href{https://github.com/OmarHassanAdelhamid/Five-of-a-Kind-capstone-project-/issues/249}{Issue \#249}
\newline\href{https://github.com/OmarHassanAdelhamid/Five-of-a-Kind-capstone-project-/issues/269}{Issue \#269}
\newline\href{https://github.com/OmarHassanAdelhamid/Five-of-a-Kind-capstone-project-/pull/442/commits/f5ab69ed317064499814f008bf446b8cde2b4949}{Commit f5ab69e}
\newline\href{https://github.com/OmarHassanAdelhamid/Five-of-a-Kind-capstone-project-/pull/442}{Pull \#442} \\ \hline
TA & MG Formatting issues (tables, vector graphics) & Addressed in pull request. & \href{https://github.com/OmarHassanAdelhamid/Five-of-a-Kind-capstone-project-/issues/272}{Issue \#272}
\newline\href{https://github.com/OmarHassanAdelhamid/Five-of-a-Kind-capstone-project-/pull/442/commits/f5ab69ed317064499814f008bf446b8cde2b4949}{Commit f5ab69e}
\newline\href{https://github.com/OmarHassanAdelhamid/Five-of-a-Kind-capstone-project-/pull/442}{Pull \#442} \\ \hline
TA & Abstract objects missing constructors & Added constructors, changed a module to a library & \href{https://github.com/OmarHassanAdelhamid/Five-of-a-Kind-capstone-project-/issues/278}{Issue \#278}
\newline\href{https://github.com/OmarHassanAdelhamid/Five-of-a-Kind-capstone-project-/pull/442/commits/9219bdc9adbb85357c92589046f6d2d218d4c8fb}{Commit 9219bdc}
\newline\href{https://github.com/OmarHassanAdelhamid/Five-of-a-Kind-capstone-project-/pull/442}{Pull \#442} \\ \hline
TA & Abstract objects need justification/explanation & Addressed in Rev 0; modules were either removed, and descriptions were added. Closed with comment. & \href{https://github.com/OmarHassanAdelhamid/Five-of-a-Kind-capstone-project-/issues/270}{Issue \#270} \\ \hline
TA & Incorrect use of environment variables & Offending module was removed in Rev 0; closed with comment. & \href{https://github.com/OmarHassanAdelhamid/Five-of-a-Kind-capstone-project-/issues/276}{Issue \#276} \\ \hline
Team & Remove unused modules & See issue; modules may be implemented further in the future. Closed with comment. & \href{https://github.com/OmarHassanAdelhamid/Five-of-a-Kind-capstone-project-/issues/431}{Issue \#431} \\ \hline
TA & Usage of ``Any''type & Removed where possible; mostly within model editing/structure related modules. & \href{https://github.com/OmarHassanAdelhamid/Five-of-a-Kind-capstone-project-/issues/428}{Issue \#428}
\newline\href{https://github.com/OmarHassanAdelhamid/Five-of-a-Kind-capstone-project-/pull/460/commits/ab0fe5610fd5f6aec7be4830c39b80b6081480b5}{Commit ab0fe56}
\newline\href{https://github.com/OmarHassanAdelhamid/Five-of-a-Kind-capstone-project-/pull/460}{Pull \#460} \\ \hline
TA & Various incorrect syntax & Fixed quotations; rest of issue addressed and closed with comments. & \href{https://github.com/OmarHassanAdelhamid/Five-of-a-Kind-capstone-project-/issues/430}{Issue \#430}
\newline\href{https://github.com/OmarHassanAdelhamid/Five-of-a-Kind-capstone-project-/pull/460/commits/ab0fe5610fd5f6aec7be4830c39b80b6081480b5}{Commit ab0fe56}
\newline\href{https://github.com/OmarHassanAdelhamid/Five-of-a-Kind-capstone-project-/pull/460}{Pull \#460} \\ \hline
Team & Include general requirements in trace tables & No modules would cover them as they apply to entire system; closed with comment. & \href{https://github.com/OmarHassanAdelhamid/Five-of-a-Kind-capstone-project-/issues/282}{Issue \#282} \\ \hline
TA & Unsure why ModelStructure, ExportStructure are ADT & Closed with comment; see comment for explanation. & \href{https://github.com/OmarHassanAdelhamid/Five-of-a-Kind-capstone-project-/issues/426}{Issue \#426} \\ \hline
TA & Unspecified datatypes & Added specification where necessary; Array3D is included in Notation. Closed with comment. & \href{https://github.com/OmarHassanAdelhamid/Five-of-a-Kind-capstone-project-/issues/429}{Issue \#429} \\ \hline

\end{longtable}

\subsection{VnV Plan and Report}

See Tables ~\ref{tab:VnVPFeedback} \&~\ref{tab:VnVRFeedback} for summaries of feedback for the VnV Plan and Report respectively.
Much of the feedback centered around making plans more specific, readable, and consistent with itself, which was taken into
consideration and implemented as suggested. See the linked issues and pull requests for more detail where necessary.

\subsubsection{VnV Plan}

\begin{longtable}{|p{2cm}|p{4cm}|p{5cm}|p{3cm}|} %------------------------------------------------------------------------
\caption{Feedback Responses and Traceability for VnV Plan}
\label{tab:VnVPFeedback} \\

\hline
\textbf{Feedback Source} & \textbf{Summary} & \textbf{Changes Made} & \textbf{Related Issues} \\ 
\hline
\endfirsthead

\hline
\multicolumn{4}{c}{\textit{Table \thetable\ continued}} \\
\hline
\textbf{Feedback Source} & \textbf{Summary} & \textbf{Changes Made} & \textbf{Related Issues} \\ 
\hline
\endhead

TA & Convert trace tables to matrices & Addressed in pull request & \href{https://github.com/OmarHassanAdelhamid/Five-of-a-Kind-capstone-project-/issues/266}{Issue \#266}
\newline\href{https://github.com/OmarHassanAdelhamid/Five-of-a-Kind-capstone-project-/pull/279/commits/6d76bf7355cf75f9e93c357bc26317ad184882ce}{Commit 6d76bf7}
\newline\href{https://github.com/OmarHassanAdelhamid/Five-of-a-Kind-capstone-project-/pull/279}{Pull \#279} \\ \hline
Peer & Add hyperlinks to trace tables & Addressed in pull request & \href{https://github.com/OmarHassanAdelhamid/Five-of-a-Kind-capstone-project-/issues/175}{Issue \#175}
\newline\href{https://github.com/OmarHassanAdelhamid/Five-of-a-Kind-capstone-project-/pull/279/commits/6d76bf7355cf75f9e93c357bc26317ad184882ce}{Commit 6d76bf7}
\newline\href{https://github.com/OmarHassanAdelhamid/Five-of-a-Kind-capstone-project-/pull/279}{Pull \#279} \\ \hline
TA & Consistent symbolic constants & Replaced all numbers throughout with symbolic constants & \href{https://github.com/OmarHassanAdelhamid/Five-of-a-Kind-capstone-project-/issues/265}{Issue \#265}
\newline\href{https://github.com/OmarHassanAdelhamid/Five-of-a-Kind-capstone-project-/pull/279/commits/6d76bf7355cf75f9e93c357bc26317ad184882ce}{Commit 6d76bf7}
\newline\href{https://github.com/OmarHassanAdelhamid/Five-of-a-Kind-capstone-project-/pull/279}{Pull \#279} \\ \hline
TA & Miscellaneous formatting issues & Implemented all mentioned in associated issue & \href{https://github.com/OmarHassanAdelhamid/Five-of-a-Kind-capstone-project-/issues/260}{Issue \#260}
\newline\href{https://github.com/OmarHassanAdelhamid/Five-of-a-Kind-capstone-project-/pull/279/commits/e17eca0087535efb0654f72a1007031ee95a9c22}{Commit e17eca0}
\newline\href{https://github.com/OmarHassanAdelhamid/Five-of-a-Kind-capstone-project-/pull/279}{Pull \#279} \\ \hline
TA & Flexibility in team roles unspecified & Added clarification to offending section & \href{https://github.com/OmarHassanAdelhamid/Five-of-a-Kind-capstone-project-/issues/261}{Issue \#261}
\newline\href{https://github.com/OmarHassanAdelhamid/Five-of-a-Kind-capstone-project-/pull/279/commits/39e45f0ffa825d9664234e5ee6788fce38648f07}{Commit 39e45f0}
\newline\href{https://github.com/OmarHassanAdelhamid/Five-of-a-Kind-capstone-project-/pull/279}{Pull \#279} \\ \hline
TA & SRS Checklist unspecific & Made more detailed, expanded upon and made explicit & \href{https://github.com/OmarHassanAdelhamid/Five-of-a-Kind-capstone-project-/issues/262}{Issue \#262}
\newline\href{https://github.com/OmarHassanAdelhamid/Five-of-a-Kind-capstone-project-/pull/279/commits/39e45f0ffa825d9664234e5ee6788fce38648f07}{Commit 39e45f0}
\newline\href{https://github.com/OmarHassanAdelhamid/Five-of-a-Kind-capstone-project-/pull/279}{Pull \#279} \\ \hline
TA & Make feedback session with supervisor more specific & Clarified, listed artifacts to elicit feedback & \href{https://github.com/OmarHassanAdelhamid/Five-of-a-Kind-capstone-project-/issues/264}{Issue \#264}
\newline\href{https://github.com/OmarHassanAdelhamid/Five-of-a-Kind-capstone-project-/pull/279/commits/39e45f0ffa825d9664234e5ee6788fce38648f07}{Commit 39e45f0}
\newline\href{https://github.com/OmarHassanAdelhamid/Five-of-a-Kind-capstone-project-/pull/279}{Pull \#279} \\ \hline
TA & Rename the Compliance Tests to a better suitable name & Rename the compliance tests to portability tests & \href{https://github.com/OmarHassanAdelhamid/Five-of-a-Kind-capstone-project-/issues/432}{Issue \#432}
\newline\href{https://github.com/OmarHassanAdelhamid/Five-of-a-Kind-capstone-project-/pull/450/changes/fa32908854b6ca3484e6380b5f4e9f774c8452ae}{Commit fa32908}
\newline\href{https://github.com/OmarHassanAdelhamid/Five-of-a-Kind-capstone-project-/pull/450}{Pull \#450} \\ \hline
Peer & VNV Readability/Formatting & Made indents + line spacing consistent throughout & \href{https://github.com/OmarHassanAdelhamid/Five-of-a-Kind-capstone-project-/issues/172}{Issue \#172}
\newline\href{https://github.com/OmarHassanAdelhamid/Five-of-a-Kind-capstone-project-/pull/460/commits/9a17d63126d54891faf5438eb9b9f26e7b860020}{Commit 9a17d63}
\newline\href{https://github.com/OmarHassanAdelhamid/Five-of-a-Kind-capstone-project-/pull/460}{Pull \#460} \\ \hline
Peer & Version tracking for traceability matrices & Will be handled in GitHub; closed with comment. & \href{https://github.com/OmarHassanAdelhamid/Five-of-a-Kind-capstone-project-/issues/174}{Issue \#174} \\ \hline
Peer & Add integration tests & Integration tests are implied by how system tests are laid out; closed with comment. & \href{https://github.com/OmarHassanAdelhamid/Five-of-a-Kind-capstone-project-/issues/173}{Issue \#173} \\ \hline


\end{longtable}

\subsubsection{VnV Report}

\begin{longtable}{|p{2cm}|p{4cm}|p{5cm}|p{3cm}|} %------------------------------------------------------------------------
\caption{Feedback Responses and Traceability for VnV Report}
\label{tab:VnVRFeedback} \\

\hline
\textbf{Feedback Source} & \textbf{Summary} & \textbf{Changes Made} & \textbf{Related Issues} \\ 
\hline
\endfirsthead

\hline
\multicolumn{4}{c}{\textit{Table \thetable\ continued}} \\
\hline
\textbf{Feedback Source} & \textbf{Summary} & \textbf{Changes Made} & \textbf{Related Issues} \\ 
\hline
\endhead

Peer & Notable spelling errors & Addressed in commit/pull request. & \href{https://github.com/OmarHassanAdelhamid/Five-of-a-Kind-capstone-project-/issues/401}{Issue \#401}
\newline\href{https://github.com/OmarHassanAdelhamid/Five-of-a-Kind-capstone-project-/pull/442/commits/9219bdc9adbb85357c92589046f6d2d218d4c8fb}{Commit 9219bdc}
\newline\href{https://github.com/OmarHassanAdelhamid/Five-of-a-Kind-capstone-project-/pull/442}{Pull \#442} \\ \hline
TA & Rename the Compliance Tests to a better suitable name & Update the VnV Report to reflect the changes made to the VnV Plan & \href{https://github.com/OmarHassanAdelhamid/Five-of-a-Kind-capstone-project-/issues/432}{Issue \#432}
\newline\href{https://github.com/OmarHassanAdelhamid/Five-of-a-Kind-capstone-project-/pull/450/changes/fa32908854b6ca3484e6380b5f4e9f774c8452ae}{Commit fa32908}
\newline\href{https://github.com/OmarHassanAdelhamid/Five-of-a-Kind-capstone-project-/pull/450}{Pull \#450} \\ \hline
TA & Assertions in VnV Report are not clear & Add more detailed explanations to the assertions in the VnV Report & \href{https://github.com/OmarHassanAdelhamid/Five-of-a-Kind-capstone-project-/issues/434}{Issue \#434}
\newline\href{https://github.com/OmarHassanAdelhamid/Five-of-a-Kind-capstone-project-/pull/450/changes/388b8b62bf0d3344265e395a87143bc13284b13e}{Commit 388b8b6}
\newline\href{https://github.com/OmarHassanAdelhamid/Five-of-a-Kind-capstone-project-/pull/450}{Pull \#450} \\ \hline

TA & Import time wording was too vague and needed better interpretation of the observed scalability trend & Revised the performance discussion to explain that import times increase sub-linearly with voxel count and clarified the interpretation of scalability behaviour. & \href{https://github.com/OmarHassanAdelhamid/Five-of-a-Kind-capstone-project-/issues/439}{Issue \#439}
\newline\href{https://github.com/OmarHassanAdelhamid/Five-of-a-Kind-capstone-project-/commit/fdf9ee2}{Commit fdf9ee2}
\newline\href{https://github.com/OmarHassanAdelhamid/Five-of-a-Kind-capstone-project-/pull/467}{Pull \#467} \\ \hline

TA & Error bars in Figures 2 and 3 were not supported by repeated trials or variance data & Removed the unsupported error bars from the figures so the graphs now show only the measured values presented in the report. & \href{https://github.com/OmarHassanAdelhamid/Five-of-a-Kind-capstone-project-/issues/438}{Issue \#438}
\newline\href{https://github.com/OmarHassanAdelhamid/Five-of-a-Kind-capstone-project-/commit/fdf9ee2}{Commit fdf9ee2}
\newline\href{https://github.com/OmarHassanAdelhamid/Five-of-a-Kind-capstone-project-/pull/467}{Pull \#467} \\ \hline

TA & Table 4 was not clearly explained and the meaning of latency was unclear & Added explicit discussion around Table 4, defined latency in this context as user interaction count rather than elapsed time, and explained what the results show about workflow efficiency. & \href{https://github.com/OmarHassanAdelhamid/Five-of-a-Kind-capstone-project-/issues/437}{Issue \#437}
\newline\href{https://github.com/OmarHassanAdelhamid/Five-of-a-Kind-capstone-project-/commit/fdf9ee2}{Commit fdf9ee2}
\newline\href{https://github.com/OmarHassanAdelhamid/Five-of-a-Kind-capstone-project-/pull/467}{Pull \#467} \\ \hline

TA & Figures should use vector graphics for better quality in the final document & Replaced the report figures with PDF vector graphics to improve clarity and rendering quality in the compiled document. & \href{https://github.com/OmarHassanAdelhamid/Five-of-a-Kind-capstone-project-/issues/436}{Issue \#436}
\newline\href{https://github.com/OmarHassanAdelhamid/Five-of-a-Kind-capstone-project-/commit/fdf9ee2}{Commit fdf9ee2}
\newline\href{https://github.com/OmarHassanAdelhamid/Five-of-a-Kind-capstone-project-/pull/467}{Pull \#467} \\ \hline

TA & Functional requirement test descriptions were too dense and some unimplemented features were incorrectly treated as executed tests & Reorganized the functional requirement test descriptions into clearer expected-result and actual-result format, and updated tests tied to unimplemented stretch-goal features to be marked as out of scope and not executed. & \href{https://github.com/OmarHassanAdelhamid/Five-of-a-Kind-capstone-project-/issues/435}{Issue \#435}
\newline\href{https://github.com/OmarHassanAdelhamid/Five-of-a-Kind-capstone-project-/commit/fdf9ee2}{Commit fdf9ee2}
\newline\href{https://github.com/OmarHassanAdelhamid/Five-of-a-Kind-capstone-project-/pull/467}{Pull \#467} \\ \hline

TA & Maintainability and security evaluation needed clearer checklist-based wording & Updated the maintainability and security sections to describe the evaluation using a structured checklist-based approach instead of only stating that the code was reviewed. & \href{https://github.com/OmarHassanAdelhamid/Five-of-a-Kind-capstone-project-/issues/433}{Issue \#433}
\newline\href{https://github.com/OmarHassanAdelhamid/Five-of-a-Kind-capstone-project-/commit/fdf9ee2}{Commit fdf9ee2}
\newline\href{https://github.com/OmarHassanAdelhamid/Five-of-a-Kind-capstone-project-/pull/467}{Pull \#467} \\ \hline
\end{longtable}

\section{Challenge Level and Extras}

\subsection{Challenge Level}
\iffalse
\plt{State the challenge level (advanced, general, basic) for your project.  Your challenge level should exactly match what is included in your problem statement.  This should be the challenge level agreed on between you and the course instructor.}
\fi

The challenge level for AutoVox is classified as \textbf{advanced}. The system integrates
multiple complex subsystems, including 3D model import and voxelization, real-
time voxel visualization and editing, magnetization and material assignment, and
structured export to a custom printer pipeline. Each subsystem involves substantial
technical and architectural considerations, such as handling large voxel datasets,
maintaining UI synchronization, and ensuring data integrity across modules. These
factors collectively establish the project as an advanced-level software system.

\subsection{Extras}
\iffalse
\plt{Summarize the extras (if any) that were tackled by this project.  Extras
can include usability testing, code walkthroughs, user documentation, formal
proof, GenderMag personas, Design Thinking, etc.  Extras should have already
been approved by the course instructor as included in your problem statement.}
\fi

Two extras were included as part of this project:
\begin{itemize}
  \item \textbf{User Manual:} A user-friendly manual that aims to act as a quick-start guide and reference for AutoVox. It includes installation instructions, explanation of use cases, and FAQs. A link to the Google Doc can be found \href{https://docs.google.com/document/d/1hXIdv2KOLxSW_-UiSePDmUuKd0TRGqq_F_Z74JUnhTE/edit?usp=sharing}{here,} and a pdf of the most recent version can also be found in docs/Extras/UserManual.
  \item \textbf{Usability Report:} A detailed report on the two usability tests conducted in the second term. It can be found on this repository under docs/Extras/UsabiltyTesting.
\end{itemize}

\section{Design Iteration (LO11 (PrototypeIterate))}
\iffalse
\plt{Explain how you arrived at your final design and implementation.  How did
the design evolve from the first version to the final version?} 

\plt{Don't just say what you changed, say why you changed it.  The needs of the
client should be part of the explanation.  For example, if you made changes in
response to usability testing, explain what the testing found and what changes
it led to.}
\fi

Our project's design went through a variety of differing versions during the timeline of 
both terms. As it will be explained, many of these changes were prompted by the team balancing
two things that conflicted at times. Our team strived to fulfull the goals that us and our 
supervisor set at the beginning of the project and tried to construct a solution that meets
every need. At the same time, our team had to consider not only technical feasibilities, but more
importantly time constraints, and what was do-able whilst managing many other workloads. The major
iterations of the project follow below.\\

\noindent\textbf{Initial Iteration \& Project Vision}
\newline
At the beginning of the project, like with any other large undertaking, our team did not have a lot 
to build off of and the size of the solution space was very large in our heads. With the added 
confusion of infrequent meetings with our supervisor due to scheduling limitations, our initial vision 
was at times misinformed. In spite of this, many core components of our project materialized early; 
the solution would be a lightweight, standalone software that allowed users to convert CAD files 
to a set of voxels that could then have magnetization vectors assigned to them. This resulted from 
preliminary discussions with our supervisor which emphasized the need for something simple, unlike the software
currently in use by their team, COMSOL. Initially, a lot of focus was placed on how CAD files would be 
voxelized, as we planned to implement this ourselves. Key ideas about the interface, such as being able 
to select layers to edit at a time, also came about from supervisor discussions, however initial designs
would centre on editing within the 3D view itself - closer to more traditional or familiar CAD workflows.\\

\noindent\textbf{Refinements after Supervisor Discussions}
\newline
Discussions with our supervisor during the development of the SRS led to a set of changes that ultimately
shaped what the design looks like today, as well as what features were included. Most notable was the 
addition of material assignment, which was left out from initial iterations as it was not mentoned in the
original project description nor brought up in conversations. As we learned more about the work being done 
and examples of models that would be built within our software, the scope began to gradually expand.
For instance, the idea of creating a project from scratch (no CAD import) resulted from conversations about 
small models with fewer than fifty voxels. The layer editing menu also resulted from conversations with one 
of our supervisor's graduate students, who likened voxel editing to being able to select cells within an excel
spreadsheet. It became clear that the interface design was key above all else, and that it needed to be 
convienient, easy to learn, and reliable. With this expansion in scope, our group made the decision to use
a pre-existing library for CAD voxelization and focus on interface iteration. In turn, this decision led to the
feature of adding and deleting voxels, an idea that paid off later in the project, anticipating the potential
for conversion errors.\\

\noindent\textbf{Scope Changes after \& during Revision 0}
\newline
During the development of revision 0 and after its presentation to the course instructors, it became clear to 
the group that the overall scope of the project would need to be reduced - not due to technical infeasibility, 
but due to time constraints. Accordingly, core functionalities we knew the interface required were focused on, 
with additional features left as potential future additions. For example, creation of projects from scratch was 
scrapped entirely as this would be a niche use case. Materials would be limited to the maximum the printer can 
currently support, and interface convienience features such as the ability to save favourite magnetizations
or resetting voxels were put on hold. This scope reduction allowed the team to focus on making the core ideas
the best they could be. These features are still retained in our documents both to maintain a history, and in the
case they will be implemented in future iterations.\\

\noindent\textbf{Changes during Usability Tests}
\newline
During the development of the VnV Report, as well as shortly before the revision 1 presentation, usability tests were
held with our supervisor and her graduate students in order to receive feedback, which are summarized in greater 
detail within the Usability Report extra. It was these meetings that led to the final implementation of AutoVox.
Overall, major changes included:

\begin{itemize}
  \item\textbf{Improvements to project creation menus.} The workflow at the time was unintuitive, and certain options 
  did not meet client needs. This was completely reworked and simplified, along with the development of a greeting screen 
  specifically targeted at new users.
  \item\textbf{Ability to rename partitions.} In order to keep track of the different parts of a project easier, instead of 
  the auto-generated positional names.
  \item\textbf{Clearer magnetization signifiers.} A checkmark was preferred at a glance, versus the original implementation 
  idea of a toggleable highlight color or directional arrow.
  \item\textbf{More guidance.} While users were able to pick up core features without instruction, it was clear that a manual
  or quick-start guide would be extremely valuable to speed up initial understanding - which wsa developed further in between tests.
\end{itemize}
\noindent\textbf{Conclusions}
\newline
Overall, AutoVox's current status reflects the balance between large goals and feasibility, incorporating valuable feedback from
stakeholders and the creativity of the team to result in a highly usable and powerful tool.

\section{Design Decisions (LO12)}
\iffalse
\plt{Reflect and justify your design decisions.  How did limitations,
 assumptions, and constraints influence your decisions?  Discuss each of these
 separately.}
\fi

Our team's primary goal when developing the system's architecture was to make future
modifications and additions as painless as possible, whilst serving the needs of our
client and meeting the technical abilities of our team. Additionally, the features
considered and how they were implemented were directly informed by the limitations and 
constraints imposed upon the project, as well as assumptions made by the team.
Below follows a brief overview of the architecture chosen and justification, followed by 
sets of limitations, assumptions, and constraints that informed feature and design choices.\\

\noindent\textbf{Architecture Overview}
\begin{itemize}
  \item\textbf{Model/View/Controller Setup.} Allows for separation of concerns between the
  user interface and voxel model updates. Given how the interface evolved considerably
  throughout the course of the project, this decision proved valuable.
  \item\textbf{Encapsulated backend structures.} Each of the primary backend components, such as
  how CAD files are processed or how project files are exported, are each broken down into
  sub-components that are encapsulated. In this way, these modules can be swapped out for alternatives
  should their performance stop meeting client needs.
\end{itemize}

\noindent\textbf{Limitations}
\begin{itemize}
  \item\textbf{Variable size of models.} This posed a two-fold issue - how to handle the files in 
  the backend efficiently, and how to effectively present large models to users without overwhelming
  them. To address this, we introduced model partitioning - breaking down a large file into a group of
  smaller ones - and the ability to zoom in or out within the layer editor.
  \item\textbf{Voxelization accuracy.} No matter how a CAD file is voxelized (eg. choice of algorithm 
  or implementation), there will be some degree of artifacting, which increases in likelihood as voxel
  size increases To address this, we included the ability to add and delete voxels. This had the additional
  benefit of allowing users to construct model supports for printing within our editor, instead of being
  forced to plan this within CAD software.
\end{itemize}

\noindent\textbf{Assumptions}
\begin{itemize}
  \item\textbf{Valid CAD file input.} AutoVox assumes STL files provided are fully enclosed and valid 
  3D structures and will not attempt to perform repairs, as how a file could be completed is up to the
  user's interpretation.
  \item\textbf{Minimum processing power requirements.} We assumed that systems running AutoVox are 
  sufficiently powerful to handle 3D-graphics, which encouraged us to adopt Three.js to implement more
  taxing model visualizations.
  \item\textbf{Library availability.} We assumed that the necessary libraries for AutoVox (eg. Trimesh)
  are readily available on target devices.
\end{itemize}

\noindent\textbf{Constraints}
\begin{itemize}
  \item\textbf{Resource usage minimization and performance considerations.} Rendering of voxels was altered
  to be done in parallel, and SQL was chosen to handle voxel models to take advantage of existing database 
  algorithms.
  \item\textbf{File stability hazards.} Our supervisor stressed that the current software in use crashed 
  frequently causing lost work - to address this, we implemented automatic saving after each edit to the 
  project. This also made saving much less taxing as saves were incremental and smaller instead of being
  accumulated over time.
  \item\textbf{Printer file intake support.} Compatibility with the existing printer software was crucial, but
  our supervisor did add that the files it could intake could be altered. To minimize the amount of alteration,
  we chose to export in a basic format that is supported by a maximal number of softwares, organized in a 
  logical and simple way.
\end{itemize}

\section{Economic Considerations (LO23)}

\subsection{Market Viability}

As technology advances, there is a growing demand for tools that streamline advanced 3D modeling workflows, particularly in emerging fields such as micro-robotics, biomedical engineering, and additive manufacturing. AutoVox targets researchers, engineers, and developers seeking to explore the expanding capabilities of 3D printing. Our software will enable the creation of innovative applications that depend on the precise control of magnetic and material properties within a structure.\\

While existing tools such as AutoCAD and CONSOL allow intuitive creation of a 3D model, they lack support for features such as automated voxelization and metadata assignment. For those constantly looking to advance in the field, experimentation is an essential part of the process when creating those advancements. The significant time required to create models during prototyping and experimentation contributes to inefficiencies in the development process and represents a direct cost to both users and organizations. To our users, AutoVox represents a way to reduce time, cost, and effort associated with model development while enabling faster experimentation. As a result, this can help streamline the usage of magnetic 3D printing workflows within different fields, increasing our software's applicability across a broader user base.

\subsection{Marketing Strategy}

To promote AutoVox, we propose:

\begin{itemize}
    \item \textbf{Academic \& Research Partnerships:} Collaborate with research labs (e.g., HeART) and universities exploring magnetization in emerging engineering domains—such as micro-robotics and advanced biomedical applications—would enable early adoption of the software within their workflows.
    \item \textbf{3D Modeling Software Partnerships:} Explore potential collaboration with 3D modeling software providers to enable seamless integration between model creation and voxel processing. While each platform addresses distinct concerns—model design and voxel-based customization—this partnership allows both tools to complement one another, promoting an efficient and cohesive workflow that benefits users of both systems.
    \item \textbf{Outreach with Physical Demonstrations:} Showcase real-world use cases (e.g., magnetic micro-robot fabrication) where the software streamlines the creation of complex components with precise and customizable magnetic properties. This may include conferences, presentations, videos, and technical reports.
    \item \textbf{Freemium Model:} Provide a free version that supports basic voxelization, allowing users to explore the software’s capabilities, while offering advanced features (e.g., metadata customization, partitioning, and export optimization) through premium tiers.
\end{itemize}    

\subsection{Production Costs}

Developing a production-ready version of AutoVox will likely requirement investment across the following areas:

\begin{itemize}
    \item \textbf{Development:} It is reasonable to assume that \$15,000–\$25,000 would be spent to extend the current prototype into a robust application. This would include UI improvements, backend optimization, and additional features. This estimate reflects freelance development rates for full-stack systems involving Python and JavaScript/TypeScript.
    \item \textbf{Maintenance \& Application Hosting:} Estimated at \$500–\$1,000 per month, covering cloud storage (for voxel datasets), backend hosting, and ongoing updates.
    \item \textbf{Research \& Testing:} Currently, the system has been tested exclusively on a 3D printer in McMaster's HeART Lab. However, for broader adoption, it must be compatible with a range of 3D printers that may vary in specifications. This means it is reasonable to assume that additional costs, estimated at around \$5,000 – \$10,000, may arise from validating the system across diverse real-world environments, including testing with different printers and simulation tools to prevent potential damage to high-cost equipment.
    \item \textbf{Marketing:} To effectively promote the product, we would focus on academic outreach, technical content creation, and in-person demonstrations. To support these efforts, an initial outreach budget of approximately \$3,000–\$5,000 is a reasonable estimate for engaging external audiences and increasing product visibility.
\end{itemize}  

Overall, first-year operational costs are estimated to be \$25,000–\$40,000, depending on feature scope and adoption scale.

\subsection{Pricing \& Profitability}

Once the software is further refined, it could adopt a freemium model to provide basic functionality at no cost while charging for advanced features and integrations. The tiers included in this model have been defined below: \\

\noindent \textbf{Free Tier:} The basic features of STL model importation, voxelization and file exportation would be included in this tier to promote accessibility and encourage adoption. This tier allows trust to be built with potential premium users by demonstrating the software’s ability to accurately process and decompose their models before they commit to purchasing the advanced features that allow them to further edit the model. \\

\noindent \textbf{Individual Tier:} Priced at \$5/month or \$50/year, this paid tier unlocks advanced features including . . .
\begin{itemize}
    \item metadata assignment for both material \& magnetization properties
    \item layer editing tools
    \item structure modification through adding and removing specific voxels
    \item a higher limit to the number of projects that can be saved at a given time in comparison to the free tier
\end{itemize}

\noindent \textbf{Enterprise \& Research Tier:}
Priced at \$500–\$1,500/year per lab or team, this tier includes:
\begin{itemize}
    \item multi-project management
    \item support for real-time collaboration, allowing multiple users within the same team to concurrently access, edit and interact with a model
    \item unlimited project creation
    \item priority support
\end{itemize}

\noindent To break even at an estimated \$40,000 annual cost, AutoVox would require:
\begin{center}
    \textit{approx. 800 users at \$50/year \hspace{0.3cm} OR \hspace{0.3cm} approx. 40 enterprise licenses at \$1{,}000/year}
\end{center}

Given the specialized but growing market in research and advanced manufacturing, this represents a realistic adoption target, particularly through academic partnerships.

\section{Reflection on Project Management (LO24)}

\subsection{How Does Your Project Management Compare to Your Development Plan}

Overall, the project followed the original Development Plan reasonably well, but with some practical adjustments. The team maintained regular meetings and communication through GitHub and messaging platforms, which aligned with the planned communication strategy. Roles were generally respected, but in practice there was flexibility, with members occasionally taking on tasks outside their original responsibilities when needed.

The planned workflow using GitHub issues, pull requests, and version control was followed consistently and became one of the strongest aspects of the project. The technologies initially selected were also used throughout the project, with no major deviations.

However, some parts of the plan, such as strict adherence to timelines and fully structured task breakdowns, were not followed as closely as intended, especially during busier academic periods.

\subsection{What Went Well?}

One of the main strengths of the project management approach was the consistent use of GitHub for tracking progress. Issues, commits, and pull requests were used effectively to organize work and maintain traceability. This made it easier to understand what changes were made and why.

Team communication was also effective, with members staying responsive and coordinating well when tasks needed to be completed quickly. Collaboration during debugging and integration phases worked particularly well, as team members were willing to support each other and resolve issues together.

Additionally, the team adapted well when requirements changed or when certain features were re-scoped, allowing the project to stay on track without major disruptions.

\subsection{What Went Wrong?}

One of the main challenges was inconsistent task planning and time estimation. Some tasks took longer than expected, while others were underestimated, which led to periods of rushed work close to deadlines.

There were also cases where features were initially planned but later identified as unnecessary or out of scope, which resulted in wasted effort early on. This was especially noticeable in the VnV process, where some tests had to be reclassified as stretch goals after development decisions changed.

Another issue was that not all documentation and testing decisions were aligned early in the process, which required revisions later to ensure consistency between the implementation and the VnV report.

\subsection{What Would you Do Differently Next Time?}

For future projects, a more realistic and flexible planning approach would be used, with better time estimation and buffer time built into the schedule. Tasks would be broken down more clearly from the beginning to reduce ambiguity and improve accountability.

There would also be earlier validation of which features are essential versus stretch goals, to avoid spending time on functionality that may later be removed or deprioritized.

In addition, stronger alignment between development and documentation would be maintained throughout the project, rather than leaving adjustments until the final stages. Regular internal reviews of documentation would help ensure consistency and reduce last-minute revisions.

Finally, maintaining the strong use of GitHub workflows would remain a priority, as it was one of the most effective aspects of the project.

\section{Reflection on Capstone}

% \plt{This question focuses on what you learned during the course of the capstone project.}

\subsection{Which Courses Were Relevant}

Several courses taken throughout the software engineering program contributed directly to the development of AutoVox.
The table below maps each relevant course to the specific aspects of the project it supported.

\begin{longtable}{|p{4cm}|p{10cm}|}
\caption{Courses Relevant to the Capstone Project}
\label{tab:RelevantCourses} \\

\hline
\textbf{Course} & \textbf{Relevance to AutoVox} \\
\hline
\endfirsthead

\hline
\multicolumn{2}{c}{\textit{Table \thetable\ continued}} \\
\hline
\textbf{Course} & \textbf{Relevance to AutoVox} \\
\hline
\endhead

SFWRENG 2AA4 -- Intro To Software Development &
Established foundational software development habits including version control with Git, modular code organization, and collaboration practices that were applied throughout the project lifecycle. \\ \hline

SFWRENG 2C03 -- Data Structures \& Algorithms &
Directly relevant to voxelization and mesh processing. Efficient spatial data structures such as hash maps were used to store and query voxel grids, and algorithmic complexity considerations informed design decisions around large 3D datasets. \\ \hline

SFWRENG 2OP3 -- Object-Oriented Programming &
The module decomposition of AutoVox into distinct, encapsulated components (e.g., ImportModule, VoxelizationModule, VisualizationModule, ExportModule) directly applied OOP principles of encapsulation, inheritance, and polymorphism covered in this course. \\ \hline

SFWRENG 3A04 -- Software Design III &
Directly applicable to the architectural design of AutoVox. Concepts from this course including module guides, module interface specifications, architectural decomposition, and information hiding were the basis for the MG and MIS documents. \\ \hline

SFWRENG 3RA3 -- Software Requirements \& Security Considerations &
Core to the SRS and Hazard Analysis deliverables. The course's framework for writing verifiable functional and non-functional requirements, identifying stakeholders, and analyzing security risks was applied directly to AutoVox's documentation. \\ \hline

SFWRENG 3S03 -- Software Testing &
Essential for the VnV Plan and VnV Report. Test design techniques such as boundary value analysis, equivalence partitioning, and unit/integration/system testing strategies were used to plan and execute the AutoVox test suite. \\ \hline

SFWRENG 4HC3 -- Human Computer Interfaces &
Directly informed the UI/UX design of AutoVox and the usability testing extra. Concepts such as Norman's design principles and cognitive load guided interface decisions and the structured usability evaluations conducted with participants. \\ \hline

SFWRENG 3DB3 -- Databases &
A crucial part of the project was the ability to store and retrieve voxel data from a database. Ensuring that the database is efficient and scalable was essential for the project. This course provided the necessary knowledge to design and implement a database for the project. \\ \hline

\end{longtable}

These courses enabled the team to be successful with the capstone project by providing the necessary knowledge and skills to complete the project. They provided a team a foundation to build upon and a framework to follow.

\subsection{Knowledge/Skills Outside of Courses}

A number of skills and areas of knowledge were required for AutoVox that were not covered in the courses taken during the program.

\begin{itemize}

  \item \textbf{3D Graphics and Rendering:} Working with a real-time 3D rendering library (Three.js) required self-directed learning. Topics such as scene graphs, camera projections, vertex and fragment shaders, instanced mesh rendering, and GPU memory management were not covered in any course and had to be learned through documentation and experimentation.

  \item \textbf{Voxelization Algorithms:} Converting arbitrary triangle meshes to volumetric voxel representations is a specialized computational geometry topic not covered in the curriculum. The product uses existing libraries and tooling for mesh-to-voxel conversion rather than a custom implementation. However, the team still had to learn concepts such as ray casting, surface-to-volume conversion, and non-smooth geometry so that requirements and design documentation could describe the system's behaviour accurately.

  \item \textbf{3D File Format Parsing:} The team did not implement an STL parser from scratch. On the server, meshes are read with the Python \texttt{trimesh} library for voxelization and dimension queries. In the client, the STL file served by the API is parsed with Three.js's \texttt{STLLoader} to build renderable geometry. Studying the STL specifications was still needed to document and reason about the file format. The team did not need to replace those libraries, but to write accurate technical documentation and interpret what the application shows.

  \item \textbf{Desktop Application Development:} Building a cross-platform desktop application using Electron required learning how to bridge web technologies (HTML/CSS/JavaScript) with native OS capabilities, manage inter-process communication (IPC) between the main and renderer processes, and handle file system access securely.

  \item \textbf{TypeScript and FastAPI:} While programming fundamentals were covered in courses, the specific TypeScript language, FastAPI framework required practical self-study to use effectively in a large multi-module project.

\end{itemize}

\end{document}